\section{Extension: research tasks and fixed automation costs}
\label{sec:task-extension}

The baseline CES treats human and automated research as two continuously variable
inputs with a constant elasticity of substitution. This is useful for characterizing
prices, shares, and balanced growth, but it suppresses an extensive margin: an AI
system may be able to perform some research tasks and not others. It also suppresses
the engineering cost required to make a task reliably automatable. This section
introduces those margins as an extension rather than changing the baseline model.
The parameter values continue to satisfy
Assumption~\ref{ass:stable-human-essential-benchmark} when the task technology is
compared with the human-essential Cobb--Douglas benchmark.
The formulation connects the paper to task-based models of automation and research
\citep{aghionjonesjones2019,acemoglu2025,jones2025rd}.

\subsection{Research tasks}

Frontier research requires a continuum of tasks indexed by $i\in[0,1]$. Let
$e_t(i)$ denote the quantity of task $i$ completed at date $t$. Effective research
is a CES aggregate across tasks,
\begin{equation}
    E_t
    =
    \left[
        \int_0^1 e_t(i)^{(\zeta-1)/\zeta}\dd i
    \right]^{\zeta/(\zeta-1)},
    \label{eq:task-research-aggregator}
\end{equation}
where $\zeta$ measures substitution across research tasks. A low $\zeta$ captures
the idea that progress can be constrained by a small set of bottleneck tasks. The
ideas production function remains
\begin{equation}
    \dot A_t=\chi A_t^\phi E_t^\eta.
    \label{eq:task-ideas-production}
\end{equation}

Humans produce task $i$ according to
\begin{equation}
    e^H_t(i)=a_H(i)h_t(i),
    \qquad
    c^H_t(i)=\frac{w_t}{a_H(i)},
    \label{eq:human-task-cost}
\end{equation}
where $c^H_t(i)$ is the human variable cost of one unit of the task. If task $i$
has been automated, research compute produces
\begin{equation}
    e^M_t(i)=a_M(i,A_t)m_t(i),
    \qquad
    c^M_t(i,A_t)=\frac{\xi_t}{a_M(i,A_t)}.
    \label{eq:machine-task-cost}
\end{equation}
Capability can increase $a_M(i,A)$, but the function is allowed to differ across
tasks. A task that the current frontier cannot perform has $a_M(i,A_t)=0$.

\subsection{Fixed costs and the automation decision}

Using AI for task $i$ requires an annuitized fixed cost $F_i(A_t)$. This cost covers
the task-specific engineering needed to construct evaluations, adapt workflows,
integrate tools, and verify reliability. For a required quantity $q_t(i)$, the two
cost schedules are
\begin{align}
    C^H_t(i,q)&=c^H_t(i)q,
    \label{eq:human-task-total-cost}\\
    C^M_t(i,q)&=F_i(A_t)+c^M_t(i,A_t)q.
    \label{eq:machine-task-total-cost}
\end{align}
The developer automates the task if
\begin{equation}
    q_t(i)\left[c^H_t(i)-c^M_t(i,A_t)\right]
    \geq F_i(A_t),
    \label{eq:task-automation-condition}
\end{equation}
provided $c^M_t(i,A_t)<c^H_t(i)$. Equivalently, the scale threshold is
\begin{equation}
    q^*_t(i)
    =
    \frac{F_i(A_t)}
    {c^H_t(i)-c^M_t(i,A_t)}.
    \label{eq:task-scale-threshold}
\end{equation}
An AI system can therefore have a lower marginal cost for a task without being used
at small scale. The fixed cost must first be amortized over enough task output. If
$F_i$ is a sunk rather than an annuitized cost, condition
\eqref{eq:task-automation-condition} compares the fixed cost with the discounted
stream of expected variable-cost savings; the economic threshold is unchanged.

Define the set of tasks automated at date $t$ as
\begin{equation}
    \mathcal A_t
    \equiv
    \left\{
        i\in[0,1]:
        c^M_t(i,A_t)<c^H_t(i),\quad
        q_t(i)\geq q^*_t(i)
    \right\},
    \qquad
    a_t\equiv\int_0^1\mathrm{I}\{i\in\mathcal A_t\}\dd i.
    \label{eq:automated-task-set}
\end{equation}
Here $\mathrm{I}\{\cdot\}$ is the indicator function. The measure $a_t$ is an
endogenous automation frontier. It expands when capability
raises task-level AI productivity, when the research wage rises relative to compute
cost, when capability lowers fixed integration costs, or when the scale of research
increases enough to amortize those costs.

\subsection{Relation to the CES benchmark}

The task model separates two margins that the baseline CES compresses into the
single share $s^M_t$. The intensive margin changes the human--machine mix within
tasks that both inputs can perform. The extensive margin changes the set
$\mathcal A_t$. When task technologies and costs have a smooth distribution, the
aggregate response of machine expenditure to $w_t/\xi_t$ can resemble the CES law
in \eqref{eq:automation-odds} over a range of the state space. The implied aggregate
elasticity need not, however, be constant: it depends on the density of tasks close
to the threshold in \eqref{eq:task-automation-condition}.

Fixed costs also explain why the initial automated share can be close to zero even
if some AI systems are technically capable of assisting research. At low research
scale, few tasks generate savings large enough to cover $F_i(A_t)$. As research
expands, more tasks cross their scale threshold. This mechanism complements the
baseline relative-price condition in
\eqref{eq:initial-human-dominance-condition}; it does not replace it.

If every task remains essential and $a_t<1$, the unautomated set can continue to
constrain effective research even when automated tasks become arbitrarily cheap.
If $a_t\to1$, the human bottleneck disappears. The asymptotic growth condition then
approaches the AI-dominated denominator $D_{\mathrm{AI}}$. If $a_t$ converges to an
interior value, the relevant feedback remains weaker and can preserve balanced
growth. Task heterogeneity therefore provides a microeconomic interpretation of
the transition between the human-dominated and AI-dominated regimes.

\subsection{Scale feedback, market structure, and equilibrium}

The fixed-cost formulation creates a scale feedback absent from the constant-
elasticity benchmark. A larger research program raises $q_t(i)$, which makes more
tasks worth automating. Automation lowers the unit cost of effective research and
can induce a further expansion of the research program. If many tasks have similar
thresholds, this feedback can produce a rapid transition or multiple local
equilibria. Whether it does is a quantitative property of the distribution of
$F_i$, $a_H(i)$, and $a_M(i,A)$; it does not follow from fixed costs alone.

The integrated monopolist may also have a private scale advantage. A task-specific
automation investment can be reused across model generations or commercial uses,
allowing the developer to spread $F_i$ over a larger base. Knowledge spillovers
create the opposite wedge: if other firms can reuse the resulting evaluations,
tools, or methods, the private developer does not capture the full return to paying
the fixed cost. The planner and the monopolist can consequently choose different
automation frontiers even at the same $A_t$ and $w_t/\xi_t$.

An equilibrium in this extension consists of the objects in Section
\ref{sec:extended}, task quantities $q_t(i)$, human and automated task inputs,
and the set $\mathcal A_t$. Conditional on aggregate states and prices, the
developer first compares the two cost schedules for every task. It then chooses the
task quantities that minimize the cost of producing $E_t$ and chooses the scale of
research using the private shadow value of capability. Market clearing and the
state equations determine the aggregate transition. A quantitative version must
calibrate the distribution of task-level fixed costs and AI productivities; the CES
model remains the parsimonious baseline until those objects can be disciplined.
